\chapter{通信工作流}

\section{普通 ask}

\begin{lstlisting}
ccb ask archi explain current architecture
\end{lstlisting}

提交后会得到 job id。可以：

\begin{lstlisting}
ccb watch <job_id>
ccb trace <job_id>
\end{lstlisting}

\section{无需结果的任务}

如果你只是让另一个 agent 做独立检查，不需要把结果带回当前任务：

\begin{lstlisting}
ccb ask --silence ops run smoke check
\end{lstlisting}

成功时静默，失败或阻塞仍应出现。

\section{需要结果的子任务}

如果当前 agent 需要另一个 agent 的结果才能继续：

\begin{lstlisting}
ccb ask --chain archi collect evidence
\end{lstlisting}

callback 不是同步等待。CCB 会记录 parent/child 关系，child 完成后把结果作为 continuation 任务送回 parent agent。

\section{Inbound task completion 和 reply routing}

当 agent 正在执行 CCB 投递的 inbound task 时，当前 provider turn 的 final
就是该 task 的 completion 内容。如果 original caller 是 registered CCB agent，
CCB 会记录 terminal reply，并按原 message lineage 投递给 caller。
Agent 不应再向 original caller 新开一个普通 \cmd{ask} 汇报完成；这样会创
建第二条任务，而不会终结当前 inbound turn。direct CLI caller 不属于这个自
动 agent-to-agent 回传范围；CLI 用户从 \cmd{watch}、\cmd{trace} 或相应
control output 读取 terminal result。

result-chain continuation 也遵循同一原则：parent agent 在 continuation 的
当前 turn 直接给出 final，由 CCB 继续向上游路由。\cmd{--chain} 用于 parent
确实需要 child 结果才能完成当前任务的场景；\cmd{--silence} 表示成功结果不需
返回，它不是向 active job 发送 correction 的通道。

当前版本没有 active-job follow-up。执行中的任务需要修正范围时，先取消旧
job，再用修正后的完整请求重新提交。不要把 correction 作为普通新任务排在同
一 agent 的 active job 后面。

job 的 \src{completed} 只表示 provider execution 正常结束。它不等同于用户
或 reviewer 已完成业务验收；业务完成度仍应根据 final 内容、产物和验证证据判
断。

CCB 在 agent materialization 时把 runtime coordination rules 写入 managed
\src{AGENTS.md}、\src{CLAUDE.md} 或相应 provider memory。更新
\src{.ccb/ccb_memory.md} 或升级模板后，先重新 materialize runtime memory；
当该 provider 不支持热加载时，再重启 provider 或启动新 session。
不要假设正在运行的 provider session 会热加载改变后的 managed memory。

\section{长内容}

长 request：

\begin{lstlisting}
ccb ask --artifact-request archi -- "$(cat long-note.txt)"
\end{lstlisting}

长 reply：

\begin{lstlisting}
ccb ask --artifact-reply archi produce detailed report
\end{lstlisting}

双向 artifact：

\begin{lstlisting}
ccb ask --artifact-io archi process this large context
\end{lstlisting}

\section{查看 queue}

查看全部：

\begin{lstlisting}
ccb queue all
\end{lstlisting}

查看某个 agent：

\begin{lstlisting}
ccb queue archi
ccb queue --detail archi
\end{lstlisting}

queue 是 mailbox-centric 视图，会叠加 runtime state 和 health。summary missing 或 error 时，视图会 degraded，可以用 \cmd{doctor storage} 继续排查。

\section{查看 inbox 和 ack}

\begin{lstlisting}
ccb inbox talk1
ccb inbox --detail talk1
\end{lstlisting}

如果 head 是已读 reply，可以 ack：

\begin{lstlisting}
ccb ack talk1
\end{lstlisting}

如果系统提示需要指定 head event：

\begin{lstlisting}
ccb ack talk1 <inbound_event_id>
\end{lstlisting}

不要试图 ack 非 head event。CCB 故意保持 head-only 语义，避免跳过更早消息。

\section{trace}

trace 支持：

\begin{lstlisting}
ccb trace sub_xxx
ccb trace msg_xxx
ccb trace att_xxx
ccb trace rep_xxx
ccb trace job_xxx
\end{lstlisting}

如果你只知道 job id，先 trace job id。trace 会显示相关 submission、message、attempt、reply、event、job。

\section{retry 和 resubmit}

当某次 attempt 失败且仍是最新 attempt：

\begin{lstlisting}
ccb retry job_xxx
\end{lstlisting}

当你要重新提交整条消息：

\begin{lstlisting}
ccb resubmit msg_xxx
\end{lstlisting}

区别：

\begin{itemize}
  \item \cmd{retry} 保留原 message lineage，增加 retry attempt。
  \item \cmd{resubmit} 创建新 message/submission/jobs。
\end{itemize}

\section{取消任务}

\begin{lstlisting}
ccb cancel job_xxx
\end{lstlisting}

取消会把 job 推进到 cancelled terminal state。它不会删除 trace lineage。

\section{推荐排障流程}

\begin{enumerate}
  \item 先 \cmd{trace <id>} 建立上下文。
  \item 如果 job 未 terminal，用 \cmd{watch <job_id>}。
  \item 如果 agent 有积压，用 \cmd{queue --detail <agent>}。
  \item 如果 caller 没看到回复，用 \cmd{inbox --detail <caller>}。
  \item 如果存储 degraded，用 \cmd{doctor storage}。
  \item 根据 trace 决定 \cmd{retry}、\cmd{resubmit}、\cmd{cancel} 或 \cmd{ack}。
\end{enumerate}
